% GENERATED FILE. Edit canonical lessons, then run npm run generate:books.
% canonical-source-hash: fnv1a64:76b31c151f34c957
% canonical-lessons: FR-C386-accuser, FR-C386-accusation, FR-C386-voleur, FR-C386-plainte, FR-C386-delit

\chapter{Les accusations}
\label{ch:fr-les-accusations}
\hlchaptermodality{386}

\begin{chapteropening}
\textbf{By the end of this chapter:} \emph{I can say who is accused of what: an accusation, a thief, a formal complaint and a minor offence.}

Complete the last lesson of chapter 386: I can say who is accused of what: an accusation, a thief, a formal complaint and a minor offence.
\end{chapteropening}

\section[accuser]{accuser — to accuse}

\label{lesson:FR-C386-accuser}



\begin{quote}
Before the new one: say the French for the certificate, then the French for the official statement.
\end{quote}

\subsection*{You'll want to know: accuser}
\textbf{accuser} — \textquotedblleft{}to accuse\textquotedblright{}.

\begin{quote}
Ils l'accusent d'avoir volé un vélo.
\end{quote}

\begin{quote}
\emph{They accuse him of stealing a bike.}
\end{quote}

\emph{Accuser quelqu'un de quelque chose}.

\begin{cousinweb}
\emph{accuser} comes through Old French \emph{acuser} (EtymDB) from Latin \emph{accūsāre} \textquotedblleft{}to call to account, to accuse\textquotedblright{}, built from \emph{ad} and \emph{causa} \textquotedblleft{}cause, lawsuit\textquotedblright{} (Lewis \& Short).
\end{cousinweb}

\begin{grammarlens}[title={What you've built}]
One more word for telling, explaining and describing at B1.
\end{grammarlens}

\subsection*{Your turn}
\begin{itemize}
\raggedright
  \item \emph{Say it:} \emph{accuser}
  \item \emph{Say it:} \emph{accuser}, once more
  \item \emph{Say it:} a sentence of your own with \emph{accuser}
  \item \emph{Recall:} say the French for the certificate, then the French for the official statement, then say \emph{accuser} again
\end{itemize}

\begin{tcolorbox}[breakable,colback=teal!4,colframe=teal!35!black,title={Before you move on}]
What does accuser mean? (\textquotedblleft{}To accuse\textquotedblright{}.) Say it once more.
\end{tcolorbox}
\section[l'accusation]{l'accusation — the accusation, the charge}

\label{lesson:FR-C386-accusation}



\begin{quote}
Before the new one: say the French for the official statement, then the French for to accuse.
\end{quote}

\subsection*{You'll want to know: l'accusation}
\textbf{l'accusation} — \textquotedblleft{}the accusation, the charge\textquotedblright{}.

\begin{quote}
C'est une accusation très grave.
\end{quote}

\begin{quote}
\emph{It is a very serious accusation.}
\end{quote}

Feminine.

\begin{cousinweb}
\emph{accusation} is a learned borrowing from Latin \emph{accūsātiō} \textquotedblleft{}an accusation, an indictment\textquotedblright{} (EtymDB; Lewis \& Short), from \emph{accūsāre} \textquotedblleft{}to accuse\textquotedblright{}.
\end{cousinweb}

\begin{grammarlens}[title={What you've built}]
One more word for telling, explaining and describing at B1.
\end{grammarlens}

\subsection*{Your turn}
\begin{itemize}
\raggedright
  \item \emph{Say it:} \emph{l'accusation}
  \item \emph{Say it:} \emph{l'accusation}, once more
  \item \emph{Say it:} a sentence of your own with \emph{l'accusation}
  \item \emph{Recall:} say the French for the official statement, then the French for to accuse, then say \emph{l'accusation} again
\end{itemize}

\begin{tcolorbox}[breakable,colback=teal!4,colframe=teal!35!black,title={Before you move on}]
What does l'accusation mean? (\textquotedblleft{}The accusation, the charge\textquotedblright{}.) Say it once more.
\end{tcolorbox}
\section[le voleur]{le voleur — the thief}

\label{lesson:FR-C386-voleur}



\begin{quote}
Before the new one: say the French for to accuse, then the French for the accusation.
\end{quote}

\subsection*{You'll want to know: le voleur}
\textbf{le voleur} — \textquotedblleft{}the thief\textquotedblright{}.

\begin{quote}
Le voleur est entré par la fenêtre.
\end{quote}

\begin{quote}
\emph{The thief came in through the window.}
\end{quote}

Feminine \emph{voleuse}. \emph{Au voleur !} — stop thief!

\begin{cousinweb}
\emph{voleur} is built from \emph{voler} \textquotedblleft{}to steal\textquotedblright{} and \emph{-eur}. \emph{voler} \textquotedblleft{}to steal\textquotedblright{} is the same word as \emph{voler} \textquotedblleft{}to fly\textquotedblright{}, from Latin \emph{volāre} \textquotedblleft{}to fly\textquotedblright{} (Lewis \& Short): the falcon that flies off with its prey.
\end{cousinweb}

\begin{grammarlens}[title={What you've built}]
One more word for telling, explaining and describing at B1.
\end{grammarlens}

\subsection*{Your turn}
\begin{itemize}
\raggedright
  \item \emph{Say it:} \emph{le voleur}
  \item \emph{Say it:} \emph{le voleur}, once more
  \item \emph{Say it:} a sentence of your own with \emph{le voleur}
  \item \emph{Recall:} say the French for to accuse, then the French for the accusation, then say \emph{le voleur} again
\end{itemize}

\begin{tcolorbox}[breakable,colback=teal!4,colframe=teal!35!black,title={Before you move on}]
What does le voleur mean? (\textquotedblleft{}The thief\textquotedblright{}.) Say it once more.
\end{tcolorbox}
\section[la plainte]{la plainte — the complaint (to the police); the moan}

\label{lesson:FR-C386-plainte}



\begin{quote}
Before the new one: say the French for the accusation, then the French for the thief.
\end{quote}

\subsection*{You'll want to know: la plainte}
\textbf{la plainte} — \textquotedblleft{}the complaint (to the police); the moan\textquotedblright{}.

\begin{quote}
Elle a porté plainte à la police.
\end{quote}

\begin{quote}
\emph{She filed a complaint with the police.}
\end{quote}

\emph{Porter plainte} — to file a complaint. From \emph{se plaindre}.

\begin{cousinweb}
\emph{plainte} comes from Old French \emph{plainte} \textquotedblleft{}lament, complaint\textquotedblright{} (EtymDB), formed on \emph{plaindre}, from Latin \emph{plangere} \textquotedblleft{}to beat the breast, to lament\textquotedblright{} (Lewis \& Short).
\end{cousinweb}

\begin{grammarlens}[title={What you've built}]
One more word for telling, explaining and describing at B1.
\end{grammarlens}

\subsection*{Your turn}
\begin{itemize}
\raggedright
  \item \emph{Say it:} \emph{la plainte}
  \item \emph{Say it:} \emph{la plainte}, once more
  \item \emph{Say it:} a sentence of your own with \emph{la plainte}
  \item \emph{Recall:} say the French for the accusation, then the French for the thief, then say \emph{la plainte} again
\end{itemize}

\begin{tcolorbox}[breakable,colback=teal!4,colframe=teal!35!black,title={Before you move on}]
What does la plainte mean? (\textquotedblleft{}The complaint (to the police); the moan\textquotedblright{}.) Say it once more.
\end{tcolorbox}
\section[le délit]{le délit — the misdemeanour, the offence}

\label{lesson:FR-C386-delit}



\begin{quote}
Before the new one: say the French for the thief, then the French for the complaint.
\end{quote}

\subsection*{You'll want to know: le délit}
\textbf{le délit} — \textquotedblleft{}the misdemeanour, the offence\textquotedblright{}.

\begin{quote}
Voler dans un magasin est un délit.
\end{quote}

\begin{quote}
\emph{Shoplifting is an offence.}
\end{quote}

Less serious than \emph{un crime}.

\begin{cousinweb}
\emph{délit} comes through Old French \emph{delit} (EtymDB) from Latin \emph{dēlictum} \textquotedblleft{}a fault, an offence\textquotedblright{}, from \emph{dēlinquere} \textquotedblleft{}to fail, to do wrong\textquotedblright{} (Lewis \& Short).
\end{cousinweb}

\begin{grammarlens}[title={What you've built}]
One more word for telling, explaining and describing at B1.
\end{grammarlens}

\subsection*{Your turn}
\begin{itemize}
\raggedright
  \item \emph{Say it:} \emph{le délit}
  \item \emph{Say it:} \emph{le délit}, once more
  \item \emph{Say it:} a sentence of your own with \emph{le délit}
  \item \emph{Recall:} say the French for the thief, then the French for the complaint, then say \emph{le délit} again
\end{itemize}

\begin{tcolorbox}[breakable,colback=teal!4,colframe=teal!35!black,title={Before you move on}]
What does le délit mean? (\textquotedblleft{}The misdemeanour, the offence\textquotedblright{}.) Say it once more.
\end{tcolorbox}
